%% ma: L10-B13-0035
%% lop: 10
%% chuong: 5
%% bai: 13
%% dang: 10.13.2
%% loai: TN4
%% muc: TH
%% dap_an: "B"
%% nguon: "chuyen-de-cac-so-dac-trung-cua-mau-so-lieu-khong-ghep-nhom-toan-10-chuong-trinh-moi.pdf (GV. Phan Nhật Linh) (P02), Bài 02 Dạng 1 Phần I Câu 25"
%% kien_thuc: "Lớp 10 Bài 13 (số trung bình)"
%% trang_thai: chinh-thuc
%% ghi_chu: "Tự duyệt theo yêu cầu giáo viên 22/10/2026."
\begin{ex}
Kết quả điểm kiểm tra môn Toán của $40$ học sinh lớp 10A được trình bày ở bảng sau:
\begin{center}
\begin{tabular}{|c|c|c|c|c|c|c|c|c|}\hline
Điểm & $4$ & $5$ & $6$ & $7$ & $8$ & $9$ & $10$ & Cộng\\\hline
Tần số & $2$ & $8$ & $7$ & $10$ & $8$ & $3$ & $2$ & $40$\\\hline
\end{tabular}
\end{center}
Số trung bình của mẫu số liệu (làm tròn đến hàng phần mười) là
\choice{$6{,}4$}{\True $6{,}8$}{$6{,}7$}{$7{,}0$}
\loigiai{$\overline x=\dfrac{4\cdot2+5\cdot8+6\cdot7+7\cdot10+8\cdot8+9\cdot3+10\cdot2}{40}=\dfrac{271}{40}=6{,}775\approx6{,}8$. Chọn B.}
\end{ex}
